\documentclass[10pt,a4paper]{amsart}
\usepackage[margin=19mm]{geometry}
\usepackage{amsmath,amssymb,mathtools}
\usepackage[scr=rsfs,cal=euler]{mathalfa}
\usepackage{xfrac}
\usepackage[unicode,hidelinks,bookmarks=false]{hyperref}
\newcommand{\ie}{i.e.\ }
\newcommand{\cf}{cf.\ }
\newcommand{\resp}{respectively\ }
\newcommand{\Subsection}{\subsection}
\providecommand{\nobreakdash}{\nobreak\hbox{-}}
\setlength{\emergencystretch}{2em}
\multlinegap0pt
\usepackage{bm}
\usepackage{bbm}
\usepackage{dsfont}
\usepackage{xspace}
\usepackage{cancel}
\usepackage{mathtools}
\usepackage{amsthm}
\makeatletter
\@ifundefined{claim}{\newtheorem{claim}{Claim}}{}
\@ifundefined{subclaim}{\newtheorem{subclaim}{Claim}}{}
\makeatother
\title[Cycle lengths in graphs of given minimum degree]{Cycle lengths in graphs of given minimum degree}
\author{Yandong Bai, Andrzej Grzesik, Binlong Li, Magdalena Prorok}
\date{Source: arXiv:2511.03085v2 (CC BY 4.0)}
\begin{document}
\maketitle

\noindent\emph{Source and licence.} Cycle lengths in graphs of given minimum degree. Version arXiv:2511.03085v2 (CC BY 4.0), distributed under the Creative Commons Attribution 4.0 International licence, as recorded in the arXiv licence metadata for this exact version and shown on the abstract page https://arxiv.org/abs/2511.03085v2. Full rights notice: licenses/arxiv-2511.03085-CC-BY-4.0.txt.\par\smallskip

\noindent\textit{Editorial scope.} Counted unit: the Claim whose source label is \texttt{claim: onecomponentH} in the article (source0.tex lines 1089--1091, source label \texttt{claim: onecomponentH}), delivered together with the article's own complete original proof of it (source0.tex lines 1093--1146, one \texttt{proof} environment of 54 lines). No mathematics is written, repaired or re-derived: statement and proof are the article's own text line for line. Required context retained uncounted: claim: no two odd paths (lines 1057--1059); claim: evendistance (lines 1071--1073); claim: mostonechord (lines 1081--1083). Every definition, display and label of those lines is retained unchanged. Omitted from this package and not redistributed: 1--1056, 1060--1070, 1074--1080, 1084--1088, 1092--1092, 1147--1594. Nothing inside a retained excerpt is omitted or altered: every retained body is delivered exactly as the article's lines are written, so no omission receipt of any kind accompanies this package. Citations: the retained excerpts carry no citation command at all, so no bibliography entry of the article is transcribed and no reference is redistributed. Reference adaptation: none was needed and none was made; every reference inside the delivered unit resolves inside it, so no unresolved reference or citation is produced. Printed slips kept literal: no printed word, symbol, hypothesis or number of the counted statement or of its proof was altered. Layout and class: the article's own class is \texttt{article} with packages including inputenc, todonotes, comment, amsmath, amsthm, amsfonts, color, amssymb, graphicx, setspace, mathrsfs, enumitem, authblk, calc; the wrapper is the lane's pinned \texttt{amsart} editorial wrapper at 10pt on A4, which restates the theorem-like environments the retained text needs and adds the article's own macro definitions that the retained text needs (none), with extra wrapper packages (bm, bbm, dsfont, xspace, cancel, mathtools). The delivered numbering is the unit's own. Assets: no retained span contains a figure environment, a graphics-inclusion command, a PDF-page inclusion, a table or a TikZ picture; no image file is copied into this package, the package declares no asset dependency and no image file is redistributed with it. Licence: version arXiv:2511.03085v2 of this article is distributed under the Creative Commons Attribution 4.0 International licence, as recorded in the arXiv licence metadata for this exact version and shown on the abstract page https://arxiv.org/abs/2511.03085v2. A full rights notice is delivered at licenses/arxiv-2511.03085-CC-BY-4.0.txt. Characters: every retained excerpt is pure ASCII, so the delivered engine, XeLaTeX with its default Latin Modern text font, reports no missing character.
\begin{claim}\label{claim: no two odd paths}
    If $C$ is an even cycle of $G$ and $P^x,P^y$ two disjoint paths from $C$ to $\{x,y\}$, say with origins $x',y'$, respectively, then $d_C(x',y')$ is even. 
\end{claim}

\begin{claim}\label{claim: evendistance}
    Every two vertices in $N_C(H)$ have an even distance in $C$.
\end{claim}

\begin{claim}\label{claim: mostonechord}
    $C$ has at most one chord.
\end{claim}

\begin{claim}\label{claim: onecomponentH}
    $G^*-V(C)$ has only one component.
\end{claim}

\begin{proof}
    Let $Q$ be a component of $G^*-V(C)$ other than $H$.

    \begin{subclaim}\label{claim: NCQgeq3}
        $|N_C(Q)|\geqslant 3$.
    \end{subclaim}

    \begin{proof}
        Assuming otherwise, by the 2-connectivity of $G$, we have $|N_C(Q)|=2$, say $N_C(Q)=\{x_1,y_1\}$. Let $G_1=G[V(Q)\cup\{x_1,y_1\}]$. Then $G_1+x_1y_1$ is 2-connected and every vertex of $G_1$ other than $x_1,y_1$ has the same degree as in $G$. It follows that $G_1$ satisfies the conclusion (1), (2) or (3). If $G_1$ contains a cycle of length 2 modulo 4, then so does $G$, a contradiction. If $G_1$ contains two $(x_1,y_1)$-paths of lengths differing by 2 modulo 4, then together with two disjoint paths from $\{x_1,y_1\}$ to $\{x,y\}$, we get two $(x,y)$-path in $G$ of lengths differing by 2 modulo 4, a contradiction. Now we conclude that $G_1$ has construction $L_i(x_1,y_1)$ for some $i=1,\ldots,4$.

        Let $C_1$ be the 4-cycle of $G_1$. If $\{x_1,y_1\}\neq\{x',y'\}$, or if $G_1$ has construction $L_i(x_1,y_1)$ with $i\in\{2,3,4\}$, then the component of $G^*-V(C_1)$ containing $w$ will contain $H$ and at least one vertex in $\{x',y'\}$, contradicting the choice of $C$. So we conclude that $\{x_1,y_1\}=\{x',y'\}$ and $G_1$ has construction $L_1(x_1,y_1)$.

        Let $G_2=G-Q$. If $x'\neq x$, then $x'$ has two neighbors in $C$ and one neighbor in $P^x$. If $y'\neq y$, then $y'$ has two neighbors in $C$ and one neighbor in $P^y$. In any case we see that every vertex of $G_2$ other than $x,y$ has degree at least 3 in $G_2$. It follows that $G_2$ satisfies the conclusion (1), (2) or (3). If $G_2$ contains a cycle of length 2 modulo 4, or $G_2$ contains two $(x,y)$-paths of lengths differing by 2 modulo 4, then so does $G$, a contradiction. Now we conclude that $G_2$ has construction $L_i(x,y)$ for some $i=1,\ldots,4$.

        Note that $C$ is an even cycle in $G_2$, so $|C|=4$ and $C$ has a chord $zz'$, where $\{z,z'\}=V(C)\backslash\{x',y'\}$. Also note that $G_1$ has construction $L_1(x',y')$, i.e., $G_1$ is a 4-cycle with a chord $z_1z'_1$, where $\{z_1,z'_1\}=V(G_1)\backslash\{x',y'\}$. Now $x'zz'y'z_1z'_1x'$ is a 6-cycle of $G$, a contradiction. 
    \end{proof}

    \begin{subclaim}\label{claim: NCQeq3}
        $|N_C(Q)|=3$ and $N_C(H)=\{x',y'\}\subset N_C(Q)$.
    \end{subclaim}

    \begin{proof}
        If $|N_C(Q)|\geqslant 4$, or $N_C(H)\backslash N_C(Q)\neq\emptyset$, then there are three vertices $u_1,u_2,u_3\in N_C(Q)$ and one vertex $z\in N_C(H)\backslash\{u_1,u_2,u_3\}$. Assume w.l.o.g. that $u_1,u_2,u_3,z$ appear in this order along $C$. It follows that there are three internally-disjoint paths $P_1,P_2,P_3$ from some $u\in V(Q)$ to $u_1,u_2,u_3$, respectively, with all internal vertices in $Q$. Now $P_1\cup P_2\cup P_3\cup\overrightarrow{C}[u_1,u_3]$ is a $\varTheta$-graph, which contains an even cycle $C_1$. The component of $G^*-C_1$ containing $w$ will contain $H$ and $z$, a contradiction. So we conclude that $|N_C(Q)|=3$ and $N_C(H)\subseteq N_C(Q)$.

        If $|N_C(H)|=3$, say $N_C(H)=\{x',y',z\}$, then $N_C(Q)=\{x',y',z\}$ as well. We can assume w.l.o.g. that $x',y',z$ appear in this order along $C$. By Claim \ref{claim: evendistance}, all three paths $\overrightarrow{C}[x',y']$, $\overrightarrow{C}[y',z]$, $\overrightarrow{C}[z,x']$ are even. Now let $P_1,P_2,P_3$ be three internally-disjoint paths from some $u\in V(Q)$ to $x',y',z$, respectively, with all internal vertices in $Q$. It follows that one of the three paths $P_1uP_2$, $P_2uP_3$, $P_3uP_1$ is even, implying that one of the three cycles $P_2uP_1x'\overrightarrow{C}[x',y']$, $P_3uP_2y'\overrightarrow{C}[y',z]$, $P_1uP_3z\overrightarrow{C}[z,x']$ is even. Let $C_1$ be such an even cycle. The component of $G^*-C_1$ containing $w$ will contain $H$ and one vertex in $\{x',y',z\}$, a contradiction. Hence we conclude that $|N_C(H)|=2$, i.e., $N_C(H)=\{x',y'\}$.
    \end{proof}

    By Claim \ref{claim: NCQeq3}, we can assume that $N_C(Q)=\{x',y',z\}$. We assume w.l.o.g. that $z\in V(\overrightarrow{C}[x',y'])$.

    \begin{subclaim}\label{claim: onecomponentQ}
        $Q$ is the only component of $G^*-V(C)$ other than $H$.
    \end{subclaim}

    \begin{proof}
        Let $Q'$ be a component of $G^*-V(C)$ other than $H$ and $Q$. By the analysis in Claims \ref{claim: NCQgeq3} and \ref{claim: NCQeq3}, $N_C(Q')=\{x',y',z'\}$ for some $z'\in V(C)\backslash\{x',y'\}$. Let $P$ be an $(x',z)$-path with all internal vertices in $Q$; and let $P'$ be an $(x',z')$-path with all internal vertices in $Q'$.

        First we assume that $z'\in V(\overrightarrow{C}[x',y'])$. We can assume w.l.o.g. that $x',z,z',y'$ appear in this order along $C$. Now $\overrightarrow{C}[x',z']\cup P\cup P'$ is a $\varTheta$-graph, which contains an even cycle $C_1$. The component of $G^*-C_1$ containing $w$ will contain $H$ and $y'$, a contradiction.

        Now we assume that $z'\in V(\overleftarrow{C}[x',y'])$. If the cycle $\overrightarrow{C}[x',z]zPx'$ is of even length, then the component of $G^*-(V(\overrightarrow{C}[x',z])\cup V(P))$ containing $w$ will contain $H$ and $y'$, contradicting the choice of $C$. So we conclude that $\overrightarrow{C}[x',z]zPx'$ is odd, implying that $Pz\overrightarrow{C}[z,y']$ is an odd $(x',y')$-path. By a similar analysis we see that $P'z'\overleftarrow{C}[z',y']$ is an odd $(x',y')$-path, contradicting Claim \ref{claim: no two odd paths} for the cycle $x'Pz\overrightarrow{C}[z,y']\overrightarrow{C}[y',z']z'P'x'$.
    \end{proof}

    \begin{subclaim}\label{claim: Ceq4}
        $|C|=4$ and $C$ has exactly one chord $zz'$, where $\{z'\}=V(C)\backslash\{x',y',z\}$.
    \end{subclaim}

    \begin{proof}
        By Claim \ref{claim: mostonechord}, $C$ has at most one chord. By Claims \ref{claim: NCQeq3} and \ref{claim: onecomponentQ}, every vertex in $V(C)\backslash\{x',y',z\}$ has no neighbors outside $C$. If $|C|\geqslant 8$, then some vertices in $V(C)\backslash\{x',y',z\}$ are not incident to a chord of $C$ and have degree 2, a contradiction. So we conclude that $|C|=4$.
        
        Set $C=x'zy'z'x'$. Notice that $z'\notin N_C(H)\cup N_C(Q)$. We have that $z'$ is incident to the chord of $C$, i.e., $zz'\in E(G)$.
    \end{proof}

    
    Let $P$ be a path from $x'$ to $z$ with all internal vertices in $Q$. Then $P\cup \{x'z,x'z',zz'\}$ is a $\varTheta$-graph, which contains an even cycle $C'$ not passing through $y'$. Consequently, the component of $G - C'$ containing $w$ is larger than $H$, which gives a contradiction.
\end{proof}

\end{document}
